\chapter{基本粒子}
\problem{11.1}{算出原子核中两个质子间的重力吸引力和静电推斥力。可以看出重力吸引力远不足以抵抗静电推斥力。这说明原子核能够稳固地结合着，必有更强的吸引力对抗库仑推斥力而有余。}
\answerarea
\problem{11.2}{在 $\pi$ 介子撞击质子的实验中，当 $\pi$ 介子的实验室能量是 $200\,\mathrm{MeV}$ 时，$\Delta$ 共振态的散射最大，求 $\Delta$ 粒子的质量。［$\pi$ 的质心坐标能量等于它的实验室能量 $\times$ 质子质量/（质子质量 $+\pi$ 的质量）。］}
\answerarea\clearpage
\problem{11.3}{在下列各式中，按照守恒定律来判断，哪些反应属于强相互作用，哪些是弱相互作用，哪些是不能实现的，并说明理由。
\begin{enumerate}
\item $\mathrm p\to\pi^++\mathrm e^++\mathrm e^-$；
\item $\Lambda^0\to\mathrm p+\mathrm e^-$；
\item $\mu^-\to\mathrm e^-+\nu_{\mathrm e}+\nu_\mu$；
\item $\mathrm n+\mathrm p\to\Sigma^++\Lambda^0$；
\item $\mathrm p+\overline{\mathrm p}\to\gamma+\gamma$；
\item $\mathrm p+\overline{\mathrm p}\to\mathrm n+\overline\Sigma{}^0+\mathrm K^0$；
\item $\mathrm K^0\to\pi^++\pi^-+\pi^0+\pi^0$。
\end{enumerate}}
\answerarea
\problem{11.4}{对表 11.2 中的每个重子，从它的有关量子数来考虑，试认出是哪三个夸克组成的。
\sourcenote{所给截图未含表 11.2，需补充或参照原书该表。}}
\answerarea\clearpage
\problem{11.5}{对表 11.2 中每个介子认出它的夸克成分。
\sourcenote{所给截图未含表 11.2，需补充或参照原书该表。}}
\answerarea
\problem{11.6}{试讨论下列三式：
\begin{enumerate}
\item $\Sigma^+\to\mathrm n+\mathrm e^++\nu_{\mathrm e}$；
\item $\Xi^-\to\mathrm n+\pi^-$；
\item $\nu_{\mathrm e}+\mathrm p\to\mathrm n+\mu^+$。
\end{enumerate}}
\answerarea\clearpage
